\documentclass{jsarticle}
\usepackage[dvipdfmx]{graphicx}
\usepackage{ascmac} % itemboxを利用するために必要

%図にEPSファイルを採用した場合，PDFファイルが生成されないことがあるので，PNGやJPEG，PDFファイルを図に採用したほうがよい．

\title{ソフトウェア設計及び実験\\前期・個人開発進捗報告書}

\author{学生番号　氏名}
\date{日付}

\begin{document}
\maketitle

\section{開発期間20xx年xx月xx日～xx日の進捗状況}

\subsection{前回からの課題と進展事項}
（記述例）前回から取り組んでいる事項は大きく分けると

\begin{itemize}
\item プレイヤキャラクタ操作とアニメーション
\item 敵キャラクタの出現と移動
\item プレイヤと敵キャラクタの衝突判定
\end{itemize}

である．
前回からの課題はすべて解決したが，新たに発生した課題は現在解決中である．
具体的な課題と進展を表\ref{table1}に示す．

\begin{table}[tb]
\begin{center}
\caption{課題と進展}
\label{table1}
\begin{tabular}{|p{7cm}|p{7cm}|}\hline
課題 & 進展 \\ \hline \hline
プレイヤをJoyConで左右移動とジャンプできるようにする & 左右移動とジャンプ，左右移動中にジャンプした時の「斜めジャンプ」も実装した．ジャンプには簡単な物理計算を取り入れた\cite{mathphys}． \\ \hline
左右移動時とジャンプ時にプレイヤキャラクタの画像を切り替えてアニメーションさせる & 左移動時には左向き，右移動時には右向きのキャラクタスプライトをそれぞれ4枚用意し，0.5sごとに手足が動いているように見せるアニメーションを実装した（図\ref{fig1}-a）．ジャンプ時は，上昇中と下降中のキャラクタ画像を用意して，切り替えるアニメーションを実装した（図\ref{fig1}-b）． \\ \hline
敵キャラをランダムな座標に描画し，キャラごとにランダムに上下左右に方向与え，0.1sごとに座標を増加または減少させて描画する & ウィンドウ内のランダムな座標に5秒おきに敵キャラを描画し移動するよう実装した（図\ref{figure2}）．0.1sごとの移動量を調整しないと，一瞬でウィンドウから消えていく．敵キャラクタごとに特徴としてスピードを変えて移動できるようにした． \\ \hline
プレイヤと敵キャラクタの中心座標を求めて，一定距離以下になると衝突したと判定する & プレイヤと敵キャラクタの画像サイズが異なるので，それぞれの画像サイズから衝突したと判定しても違和感のない距離を手動で設定し，if文で判定することで衝突判定を実装した． \\ \hline
敵キャラが増える衝突判定の処理が重くなる課題が新たに発生した & プレイヤキャラクタの近く（±50ピクセル）にいる敵キャラクタだけに対して衝突判定しようとしているが，すべての敵キャラとの座標は比較しているので，処理速度はまだ遅い．別の方法が必要だと思われる． \\ \hline
\end{tabular}
\end{center}
\end{table}

\begin{figure}[t]
\begin{center}
\includegraphics[width=230pt]{player.png}
\caption{プレイヤキャラクタの移動　(a) 左右移動時，(b) ジャンプ時}
\label{fig1}
\end{center}
\end{figure}

\begin{figure}[t]
\begin{center}
\includegraphics[width=230pt]{emenies.png}
\caption{描画された敵キャラクタ（起動度約1分後の画面）}
\label{fig2}
\end{center}
\end{figure}

\clearpage

\begin{itembox}[l]{衝突アルゴリズム}
\begin{verbatim}
for(int i = 0; i < enemy_num; i++){
...
}
\end{verbatim}
\end{itembox}



\subsection{特筆すべき事項}
（記述例）このゲームは敵キャラを避けまくってゴールにたどり着くことに面白さを置いているため，プレイヤキャラクタのスムーズかつ直観的な移動は大前提であり，これを達成できたと考えている．

本開発期間では作業が比較的スムーズに進んだため，新たな課題の解決と同時並行で，ガントチャートで次週に取り組む事項（以下に列挙）にも取り組んだ．未完成ではあるが，今週の前半には完成する見込みである．

\begin{itemize}
\item 衝突判定による画面の視覚的演出（衝突時に一瞬ウィンドウ一面を赤で描画）とHPゲージの更新描画
\item 衝突時の効果音再生
\item プレイ時間のカウントダウンとその描画
\end{itemize}


\section{開発期間20xx年xx月xx日～xx日の予定}

\subsection{完了させる事項}
（記述例）ゲームの調整に最終週を充てたいので，この週は以下の事項を必ず完了させる．

\begin{itemize}
\item 衝突判定による画面の視覚的演出（衝突時に一瞬ウィンドウ一面を赤で描画）とHPゲージの更新描画
\item 衝突時の効果音再生
\item プレイ時間のカウントダウンとその描画
\item 終了（勝利/敗北）判定と終了画面の表示（その後，スタート画面に戻る）
\end{itemize}

\subsection{継続して取り組む事項}
（記述例）キャラクタの画像は自作だが，もう少し見た目を改善したいので，以下に示すプログラムと同時並行へ画像作成にも取り組む．

\begin{itemize}
\item 敵キャラクタごとの衝突時ダメージの設定
\item 敵キャラクタの移動パタンの増加（幾何学的な模様のように動かしたい）
\item プレイヤのランキング表示
\end{itemize}

\subsection{ガントチャート}
（記述例）新たな課題の発生はあったが，スムーズに開発が進んだので，以前のガントチャートよりも，一通り完成の日が2，3日ほど早くなる予定に変更している（図\ref{fig3}）．しかし，思わぬバグが入り込んでバグ取りに時間がとられる可能性もあり，ゲームバランスの調整にも時間をかけたいので，現在よりもいくらか早いペースで開発を進めたい．

\begin{figure}[t]
\begin{center}
\includegraphics[width=230pt]{ganttchart.png}
\caption{更新したガントチャート}
\label{fig3}
\end{center}
\end{figure}


\begin{thebibliography}{9}
\bibitem{mathphys} 堂前 嘉樹 （著）: ゲームを動かす数学・物理 R,  ボーンデジタル (2020).
\end{thebibliography}

\end{document}

